\begin{proof}
\begin{enumerate}[label=\arabic*.,leftmargin=*]
\item The base is intrinsic for these sphere members. A biholomorphism $F:X\to X'$ gives $g=f'F$. Since $H^2(S^6,\mathbb R)=0$, the closed form $g^*\omega_{\rm FS}$ is exact. Restrict to a regular torus $T$ and wedge with a flat K\"ahler form on $T$. Stokes' theorem gives integral zero, while the integrand is nonnegative and strictly positive wherever $dg|_T\ne0$. Thus $g$ is constant on each regular fibre. The image of $(f,g)$ is a compact analytic curve whose projection to $B$ has degree one; its normalization gives a holomorphic $h:B\to B$ with $g=hf$. Apply the same construction to $F^{-1}$ to see that $h$ is an automorphism. The three homological local monodromy orders are three, four and infinity, independently of $k,l$; their intertwining maps under $F$ therefore fix each special point. Hence $h=1$. The nonreduced fibre multiplicities now force equal $k,l$. This is the base-preservation argument of \cite[Proposition~4.7, pp.~14--15]{ViaclovskyLog2026}, with the displayed maps.

\item The root model of degree $m_1$ is the pullback of the degree-three model by $s\mapsto s^k$, and likewise at the other point. Its linear action remains $A_1$ or $A_2$, while its root generator has order $m_j$. The primitive invariant coordinate $r_j/m_j$ proves freeness of every nonidentity power. The normalization and deck-commutation arguments of \ref{fam:src:A.1.2} and \ref{fam:src:A.2.1} apply to these root covers without changing a generator. Their norms are
\[
 N_{3k}=kN_1,\qquad N_{4l}=lN_2,
\]
so the fixed-coordinate norm images are $3k\mathbb Z\varepsilon+k\mathbb Zd$ and $4l\mathbb Z\varepsilon'+2l\mathbb Zd$. The norm equation (B2.1), with $m_j$ replaced by these covering degrees, gives exactly the displayed congruences. On ordinary fibres the affine lift has simultaneous rational centralizer
$\mathbb Q I+\mathbb Q E$, as in \cref{fam:src:A.2.1}.
Comparison of the constant period columns gives fibre-linear part $aI$,
integer shift $n$ and lattice map $aS_n^{-1}$, where $S_n=I+nE$.
Functorial normalization of the root base change makes the lift commute
with the prescribed deck generator; summing its defect gives the norm
equation. The explicit base norms are
\[
 N_3v=3\gamma(v)\varepsilon+(2v_u+v_w+3v_\delta)d,\qquad
 N_4v=4\gamma(v)\varepsilon'+(2v_u+2v_w+4v_\delta)d,
\]
where $N_3=N_1$ and $N_4=N_2$ in the notation above.
At the cusp the character ratio of any extending map is $t_c^q$ times
a holomorphic unit, by \cref{fam:src:B.2.1}; its winding vanishes in
the quotient by the toric directions. The middle cocycle then forces
$p'=ap$.

\item The integer $p$ records the primitive sphere relation in this marking. The common torus coinvariants are detected by $\gamma$, since $(A_1-I)\Lambda+(A_2-I)\Lambda=\ker\gamma$ by their explicit saturated bases. The finite meridian relations are $x^{m_1}=c^{r_1}$, $y^{m_2}=c^{r_2}$, and the cusp relation is $xy=c^{l_0}$ with $c$ central. Eliminating $y$ gives an abelian relation matrix with determinant $-p$ and primitive first elementary divisor because $\gcd(m_1,m_2)=1$. Thus the sphere data have $p=\pm1$. Reducing $p'=ap$ modulo $m_1,m_2$ gives the two congruences for $r_j$.

\item For the local sufficiency maps put $\alpha=a$, $h=n$,
$S_h=I+hE$ and $d_j=S_hw'_j-\alpha w_j$.
Suppose $p'=\alpha p$ and
$hr'_1\equiv\alpha s_1-s'_1\pmod k$,
$hr'_2\equiv\alpha s_2-s'_2\pmod{2l}$.
Reduction of $p'=\alpha p$ modulo $m_1,m_2$ gives
$r'_j\equiv\alpha r_j\pmod{m_j}$.
We may choose integral solutions of $N_{m_j}\ell_j=-d_j$ explicitly:
\[
 \begin{split}
 \ell_1&=\left(\frac{\alpha r_1-r'_1}{m_1},0,
 -\frac{s'_1+hr'_1-\alpha s_1}{k},0\right)^t,\\
 \ell_2&=\left(\frac{\alpha r_2-r'_2}{m_2},
 -\frac{s'_2+hr'_2-\alpha s_2}{2l},0,0\right)^t.
 \end{split}
\]
Put $y_j=d_j/m_j+\ell_j$; these vectors belong to the rational
kernel of $N_{m_j}$ and have first coordinate zero.
The following rational vectors solve $(I-A_j)b_j=y_j$:
\[
 \begin{split}
 b_1&=(0,(2y_{1u}+y_{1w})/3,(y_{1w}-y_{1u})/3,0)^t,\\
 b_2&=(0,(y_{2u}-y_{2w})/2,(y_{2u}+y_{2w})/2,0)^t.
 \end{split}
\]
The last coordinates of the equations follow from the displayed
norm-kernel equations.  The maps
\[
 (s,\zeta)\longmapsto(s,\alpha\zeta+\Pi_c(s)b_j)
\]
are holomorphic maps of the normalized torus families over the
degree-$m_j$ root discs.  Their inverses are holomorphic, and period
covariance and $(I-A_j)b_j=d_j/m_j+\ell_j$ give exact commutation
with the prescribed affine deck generators, modulo the integral
period $\ell_j$.  They therefore descend to biholomorphisms of the
finite fillings, including the full multiple fibres.

\item On the punctured root collar the corresponding middle translation is
\[
 \mathcal A_j=\Pi_cb_j+\frac{\log s_j}{2\pi i}\Pi_cd_j.
\]
Its integral defect is $\ell_j$.  The middle cocycle relation fixes
$\ell_0=-A_0(\ell_1+A_1\ell_2)$.  In the quotient
$\overline\Lambda=\Lambda/(\mathbb Z\widehat w+\mathbb Z\widehat\delta)$,
the first-coordinate condition is exactly $p'=\alpha p$.
Adding an integral multiple of $(0,1,-2,0)^t\in\ker N_{m_1}$ to
$\ell_1$ changes the other quotient coordinate by a generator,
without changing this first coordinate.  We can consequently arrange
\[
 [\ell_0]=-\alpha c(\sigma)+c(\sigma')\quad\hbox{in }\overline\Lambda.
\]
An integral toric vector then supplies $q\in\mathbb Z^2$ such that
the cusp reference $sq+\alpha\sigma-\sigma'$ has precisely defect
$\ell_0$.  Here $s=(2\pi i)^{-1}\log t_c$ is the cusp logarithm;
the maps of cusp toric covers are
\[
 (x,t_c)\longmapsto(t_c^q u(t_c)x^\alpha,t_c),
 \qquad u(0)\in(\mathbb C^*)^2.
\]
Integer translations of the cocharacter lattice, its inversion,
and multiplication by a holomorphic unit preserve the specified
hexagonal fan and all its strata.  These maps and their inverses
extend holomorphically across the three double curves and both
triple points.

\item Differences of allowed
finite extensions are the equivariant holomorphic vertical vectors
on the complete root discs.  The extra subgroup of order $k$ at
the first point acts trivially on the vertical vector coefficients
and rotates $s$; its invariant holomorphic coefficients are exactly
those in $s^k$.  The remaining order-three equivariance is the
original one.  The argument at the second point uses $s^l$ and
order four.  These invariant-coefficient maps identify the coarse difference sheaf
with $\mathcal V_B$ for multiplicities $(3,4)$.  Its first-coordinate quotient is
$\mathcal O_B(-p_0)$: the modular frame has vanishing orders $2k,l$
on the two enlarged roots and a simple cusp pole.  Its fixed
second-coordinate subline is $\mathcal O_B$.  Finite averaging
provides local lifts, hence
\[
 0\longrightarrow\mathcal O_B\longrightarrow\mathcal V_B
 \longrightarrow\mathcal O_B(-p_0)\longrightarrow0.
\]
Both end terms have vanishing $H^1$ on $\mathbb P^1$, so
$H^1(B,\mathcal V_B)=0$.  The permitted local translations form a
torsor under this sheaf: on each overlap their defects agree, so
their difference is a single-valued section of the indicated
difference sheaf.  Vanishing of its \v Cech class gives a global
middle translation.  Its overlap equalities are exactly the two
finite formulas and the cusp formula just displayed.  Consequently
the local biholomorphisms glue, with their inverses, over $B$.

\item For identical data, $p\ne0$ excludes $a=-1$. The congruences reduce to $k\mid nr_1$ and $2l\mid nr_2$. The two primitivity assumptions imply $k\mid n$ and $2l\mid n$. Since $\gcd(k,2l)=1$, their least positive common multiple is $2kl$. For different sphere translation data, $r'_1$ is invertible modulo $k$, $r'_2$ modulo $2l$, and the Chinese remainder theorem gives exactly one integer shift class modulo $2kl$. The values of $s_1$ modulo $k$ and $s_2$ modulo $2l$ all occur without affecting freeness or $p$. Thus there are exactly $2kl$ classes at fixed $u$. The change $c\mapsto c-1$ on a positive $u$ circuit acts by a generator of this cyclic residue group.

\end{enumerate}
\end{proof}